\begin{abstract}
Biological collections hold hundreds of millions of specimens whose label
metadata remain largely ``dark data.'' At the Komarov Botanical Institute
(LE), the world's fifth-largest herbarium, fewer than 1\,\% of the mycological
holdings are imaged and only $\sim$22\,\% carry structured metadata. We present
a systematic study of automatic transcription and structuring of herbarium
label metadata on a demanding real-world corpus: 394 type specimens of smut
fungi (Ustilaginomycotina) from LE, spanning 1873--2009, four languages
(Latin, German, English, and Russian including pre-reform orthography), and
four text types (typeset, typewritten, handwritten, mixed). Across
\textbf{17 experiments} we compare (i) a seven-engine classical OCR ensemble
with orientation preprocessing, (ii) 18 vision-language models (VLMs) as
normalizers in zero-shot, few-shot, and OCR-fusion regimes, (iii) two-stage
OCR-VLM+formatter pipelines, and (iv) parameter-efficient low-rank adaptation
(LoRA) of Qwen3-VL-8B-Instruct selected by an 83-configuration sweep. All
systems are scored against a 21-field Darwin~Core-aligned schema with a single
unified metric suite. Classical OCR reaches at best $F_1=0.767$; a zero-shot
VLM $0.807$; five-shot prompting adds $+6.7$ points to $0.874$; and a LoRA
adapter trained on only 310 examples in $\sim$25\,min on one GPU attains
$F_1=\textbf{0.972}$ ($\mathrm{CER}=0.061$) on the full corpus---a $+24.0$-point
gain over the OCR baseline, beating 18 stronger or larger models. We show that
merging the adapter and batching inference yields a $5\times$ speedup
($2.77$\,s/image) at negligible quality loss; that a GBIF taxonomic validator
is a quality \emph{signal} rather than an $F_1$ booster; and that a naive LLM
corrector \emph{hurts} by substituting modern synonyms for historical names.
A four-model consensus audit surfaces 51 likely ground-truth omissions.
Finally, applying the production model to 534 unseen \emph{Anthracoidea}
specimens yields a $99.8\,\%$ parse rate with per-field behavior consistent
with the annotated corpus. Our results provide a practical, on-premise recipe
for cost-effective mobilization of mycological label data.
\end{abstract}
